% PLACEHOLDERS for Chapters 1 to 5 (pilot). Each holds only the statements that
% Chapter 6 cites, copied from the lecture notes with their labels. References to
% material that is not in the book yet were removed from the copies. When a chapter
% is written, its placeholder goes and the statements keep their labels.

% ---------------------------------------------------------------------------
\chapter{Rings, Homomorphisms and Ideals}\label{ch:rings}

\begin{stubnote}
\textbf{Placeholder.} To be written from Week~1 Wednesday and Week~2 Tuesday. Only the statements Chapter~6 refers to are copied here.
\end{stubnote}

\begin{remark}[Standing Convention for Rings in this Course]\label[remark]{rem:ring-convention}
In the general mathematical literature, rings are not necessarily required to have commutative multiplication, nor are they always required to possess a multiplicative identity $1$. In this course, unless explicitly stated otherwise, the term \newterm{ring} will always designate a \textbf{commutative ring with unit}: a ring with unit $1$ such that $a \cdot b = b \cdot a$ for all $a, b \in R$. The zero ring ($1 = 0$, \cref{prop:zero-ring}) is still allowed.
\end{remark}

\begin{proposition}[Characterization of the Zero Ring]\label[proposition]{prop:zero-ring}
Let $R$ be a ring. If $1 = 0$ in $R$, then $R$ consists of a single element:
\[
R = \{0\}.
\]
\end{proposition}

\begin{definition}[Field]\label[definition]{def:field}
A \newterm{field} is a commutative ring $(R, +, \cdot, 0, 1)$ with $0 \neq 1$ in which every non-zero element possesses a multiplicative inverse:
\[
\forall a \in R \setminus \{0\} \; \exists a^{-1} \in R \colon a \cdot a^{-1} = 1.
\]
\end{definition}

\begin{theorem}[Ideals in a Field]\label[theorem]{thm:ideals-in-a-field}
Let $F$ be a field. The only ideals of $F$ are the zero ideal $\{0\}$ and the unit ideal $F$:
\[
I \subset F \text{ an ideal} \implies I = \{0\} \quad \text{or} \quad I = F.
\]
\end{theorem}

\begin{definition}[Principal Ideal Generated by an Integer]\label[definition]{def:principal-ideals-z}
For any integer $n \in \mathbb{Z}$, the set of all multiples of $n$ is denoted:
\[
n\mathbb{Z} := \{ n \cdot k \mid k \in \mathbb{Z} \} \subset \mathbb{Z}.
\]
\end{definition}

\begin{theorem}[Ideals of $\mathbb{Z}$ are Principal]\label[theorem]{thm:ideals-of-z-are-principal}
Every ideal $I \subset \mathbb{Z}$ is of the form $I = n\mathbb{Z}$ for a unique integer $n \ge 0$.
\end{theorem}

\begin{definition}[Compactly Supported Smooth Functions]\label[definition]{def:compact-support-ideal}
The \newterm{support} of a function $f \colon \mathbb{R} \to \mathbb{R}$ is the closure of the set where $f$ does not vanish:
\[
\operatorname{supp}(f) := \overline{\{ x \in \mathbb{R} \mid f(x) \neq 0 \}}.
\]
The set of smooth functions with compact support is denoted:
\[
I := \{ f \in C^\infty(\mathbb{R}) \mid \operatorname{supp}(f) \text{ is compact in } \mathbb{R} \}.
\]
By Heine--Borel, $\operatorname{supp}(f)$ is compact if and only if it is bounded.
\end{definition}

\begin{proposition}[$I$ is an Ideal, but not a Subring]\label[proposition]{prop:compact-support-ideal}
The set $I \subset C^\infty(\mathbb{R})$ of compactly supported smooth functions is an ideal of $C^\infty(\mathbb{R})$, but is not a subring.
\end{proposition}

% ---------------------------------------------------------------------------
\chapter{The Language of Categories}\label{ch:categories}

\begin{stubnote}
\textbf{Placeholder.} To be written from Week~2 Tuesday and Wednesday.
\end{stubnote}

% ---------------------------------------------------------------------------
\chapter{Integral Domains, Units and Divisibility}\label{ch:domains}

\begin{stubnote}
\textbf{Placeholder.} To be written from Week~2 Wednesday and Week~3 Tuesday. Only the statement Chapter~6 refers to is copied here.
\end{stubnote}

\begin{definition}[Integral Domain]\label[definition]{def:integral-domain}
A commutative ring $R$ is an \newterm{integral domain} (or \newterm{domain}) if:
\begin{enumerate}[label=\textbf{(\roman*)}]
    \item $R \neq 0$ (that is, $1_R \neq 0_R$: the zero ring is excluded; cf.\ \cref{prop:zero-ring}).
    \item For all $a, b \in R$:
    \[
    a \neq 0 \quad \text{and} \quad b \neq 0 \implies a \cdot b \neq 0.
    \]
\end{enumerate}
\end{definition}

% ---------------------------------------------------------------------------
\chapter{Constructing Rings: Fractions, Polynomials, Power Series}\label{ch:constructions}

\begin{stubnote}
\textbf{Placeholder.} To be written from Week~3 Tuesday and Wednesday and Week~4 Tuesday.
\end{stubnote}

% ---------------------------------------------------------------------------
\chapter{Quotient Rings and the Correspondence Theorem}\label{ch:quotients}

\begin{stubnote}
\textbf{Placeholder.} To be written from Week~4 Tuesday and Wednesday and Week~5 Tuesday. Only the statements Chapter~6 refers to are copied here.
\end{stubnote}

\begin{definition}[Ideal Generated by Finitely Many Elements]\label[definition]{def:ideal-generated-by-elements}
Let $x_1, \dots, x_n \in R$. The \newterm{ideal generated by} $x_1, \dots, x_n$, denoted by $(x_1, \dots, x_n)$, is the set of all $R$-linear combinations of the generators:
\[
(x_1, \dots, x_n) := \left\{ \sum_{i=1}^n a_i x_i \;\middle|\; a_1, \dots, a_n \in R \right\}.
\]
When $n = 1$, the ideal $(x) = \{ a x \mid a \in R \} = x R$ is called a \newterm{principal ideal} (for $R = \mathbb{Z}$, these are the ideals $n\mathbb{Z}$ of \cref{def:principal-ideals-z}).
\end{definition}

\begin{theorem}[The Quotient Ring \texorpdfstring{$(\faktor{R}{I}, +, \cdot, \bar{0}, \bar{1})$}{(R/I, +, ., 0, 1)}]\label[theorem]{thm:quotient-ring-structure}
Let $R$ be a commutative ring with identity, and let $I \subset R$ be an ideal. The quintuple $(\faktor{R}{I}, +, \cdot, \bar{0}, \bar{1})$, equipped with the operations on classes
\[
\bar{x} + \bar{y} \stackrel{\text{def}}{=} \overline{x + y}, \qquad \bar{x} \cdot \bar{y} \stackrel{\text{def}}{=} \overline{x \cdot y}, \qquad 0_{\faktor{R}{I}} := \bar{0}, \qquad 1_{\faktor{R}{I}} := \bar{1},
\]
forms a commutative ring with identity.
\end{theorem}

\begin{definition}[Quotient Map]\label[definition]{def:canonical-quotient-map}
Let $R$ be a ring and $I \subset R$ an ideal. The \newterm{quotient map} is defined by:
\[
q \colon R \to \faktor{R}{I}, \quad x \mapsto q(x) := \bar{x} = x + I.
\]
\end{definition}

\begin{notation}[Sets of Ideals]\label[notation]{not:ideal-sets-k}
Let $R$ be a ring. The set of all ideals of $R$ is denoted by $\mathscr{I}(R)$. For an ideal $K \subset R$, the ideals of $R$ containing $K$ form the set
\[
\mathscr{I}_K(R) := \{ I \subset R \mid K \subset I \text{ and } I \text{ is an ideal of } R \}.
\]
\end{notation}

\begin{theorem}[Ideal Correspondence Theorem]\label[theorem]{thm:ideal-correspondence}\label[theorem]{prop:correspondence-theorem-ideals}
Let $R$ be a ring and $K \subset R$ an ideal, with quotient map $q \colon R \to \faktor{R}{K}$. The assignment
\[
\begin{aligned}
\mathscr{I}_K(R) &\longleftrightarrow \mathscr{I}\left(\faktor{R}{K}\right) \\
I &\longmapsto q(I) = \faktor{I}{K} \\
q^{-1}(J) &\mathrel{\reflectbox{$\longmapsto$}} J
\end{aligned}
\]
is an inclusion-preserving bijection between the set of ideals of $R$ containing $K$ and the set of all ideals of the quotient ring $\faktor{R}{K}$.
\end{theorem}
